\section{Introduction}
\label{sec:intro}
Immersive virtual reality (VR) has moved from entertainment into operational, safety-critical settings such as clinical rehabilitation, industrial and defence training and human-performance assessment, where the headset actively shapes perception, balance and spatial orientation rather than merely presenting information on a screen. This makes VR vulnerable in a way conventional software is not: an adversary who manipulates rendering parameters such as frame pacing, field of view or camera motion can deliberately induce cybersickness---nausea, disorientation, oculomotor strain and, in severe cases, postural instability---and thereby cause physiological harm~\cite{tseng2022dark,odeleye2021framerate,kundu2025securing}. Long studied as an ergonomic side effect of sensory conflict~\cite{caserman2021cybersickness}, cybersickness thus also becomes a perception-layer attack surface~\cite{odeleye2023virtually,giaretta2025security}. We use the term \emph{adversarial cybersickness} for the deliberate induction, amplification or exploitation of these symptoms.

AI-based detection has advanced from questionnaires and classical classifiers to deep temporal and multimodal models~\cite{dennison2016physio,islam2020automatic,islam2021fusion,liu2024cnneca,yalcin2024automatic,jeong2024precyse}. Three limitations persist. First, most detectors are static classifiers whose evaluation protocol is rarely participant-independent~\cite{yang2022mlreview}. Second, they are evaluated under non-adversarial assumptions, so high reported accuracy does not imply robustness to deliberate manipulation~\cite{kundu2025securing}. Third, the costs of the two error types are treated as equal, although a missed cybersickness event leaves a user exposed to harm whereas a false alarm causes a brief, recoverable interruption.

These gaps motivated E-DDQN, a safety-prioritised dueling Double DQN~\cite{vanhasselt2016double,wang2016dueling} trained with an asymmetric reward that penalises missed events five times more than false alarms, with class-balanced replay, soft target updates, validation-based selection and post hoc SHAP explanations. An earlier version of this work reported 96.88\,\% accuracy and 99.53\,\% recall for E-DDQN and an advantage over a standard DQN. Re-examining that study raised five questions: whether the formulation is genuinely reinforcement learning, whether the evaluation leaked information across sessions and participants, whether the test set was used during development, whether the baseline comparison was controlled, and whether the adversarial framing and several citations were supported. This paper answers them by rebuilding the study from the raw data.

The contributions are:
\begin{enumerate}
\item A safety-prioritised formulation of cybersickness detection under an adversarial threat model, together with a proof that, on logged data, the alarm does not affect future observations, so the problem is a contextual bandit, and that the 5:1 asymmetric reward induces exactly the Bayes threshold $P(y=1\mid s)\ge1/6$ of 5:1 cost-sensitive classification (Section~\ref{sec:formulation}).
\item A corrected E-DDQN (within-session transitions with terminal flags, removal of an invalid teacher-forcing scheme, full-information replay as its principled replacement, correct target-network normalisation statistics), verified by executable unit tests (Section~\ref{sec:eddqn}).
\item A leakage-free, participant-conservative evaluation with a protocol frozen before any model was trained on the evaluation data on the public UFFCSData dataset---session-linkage groups, nested grouped cross-validation, causal features and train-only preprocessing---comparing E-DDQN with a standard DQN, a contextual-bandit variant and cost-weighted supervised models with identical tuning budgets, operating rules and paired statistics (Sections~\ref{sec:groups}--\ref{sec:protocol}).
\item Reproducible, bounded telemetry-perturbation stress tests of the detectors (frame pacing, camera motion, locomotion), including adaptive evasion and false-alarm worst cases (Section~\ref{sec:threat}).
\item An audit showing that the previously reported performance is explained by oversampling before a row-level split, and a SHAP analysis of the decision margin rather than of a class score.
\end{enumerate}

Under the leakage-free protocol the problem is hard: the best ROC-AUC is \HGBCostRocAucMean{}, and only a cost-weighted random forest improves materially on always alarming under the 5:1 cost (cost \RFCostCostMean{} vs.\ \AlwaysCost; 95\,\% CI \RFCICostLo--\RFCICostHi). E-DDQN reaches cost \EDDQNCostCostMean{} and ROC-AUC \EDDQNCostRocAucMean{}, does not outperform the standard DQN, and is significantly worse than the random forest. We report these negative results because they are what the evidence shows. Section~\ref{sec:related} reviews related work, Section~\ref{sec:method} presents the methodology, Section~\ref{sec:setup} the setup, Section~\ref{sec:results} the results and Section~\ref{sec:discussion} the discussion; Section~\ref{sec:conclusion} concludes.
